\begin{abcliste}
\abc Die vorbeiströmende Luft nimmt bei konstantem Druck pro Sekunde die Abwärme $P$ auf. Mit $c_p = \cpLO$ (FoTaBe S.~191) ist der nötige Massenstrom
\begin{align}
 \dot m &= \mfF = \frac{\PA}{\cpL\cdot\left(\TbO-\TaO\right)} = \mf
\end{align}
Die Dichte der Luft ist $\rho_0 = \rLoO$ bei $T_0 = \TnO$ und Normaldruck (FoTaBe S.~187). Bei konstantem Druck ist die Dichte umgekehrt proportional zur absoluten Temperatur, für die einströmende Luft mit $T_1 = T_0+\vartheta_1 = \TaK$ also
\begin{align}
 \rho_1 &= \rLaF = \rLo\cdot\frac{\Tn}{\TaK} = \rLa
\end{align}
Der Volumenstrom ist damit
\begin{align}
 \dot V &= \VfF = \frac{\mf}{\rLa} = \Vf \approx \result{\VflP{0}{2}}
\end{align}
Pro Sekunde müssen also etwa \VflP{0}{2} Luft am Motor vorbeiströmen.

\abc Die Luft strömt mit der Geschwindigkeit $v$ durch die Kühlerfläche $A$, also ist $\dot V = A\,v$ und
\begin{align}
 A &= \AKF = \frac{\Vf}{\vL} = \AK \approx \result{\AKcP{0}{2}}
\end{align}

\abc Ja. Bei höherer Fahrgeschwindigkeit strömt pro Sekunde mehr Luft durch den Kühler ($\dot V = A\,v$). Die gleiche Abwärme erwärmt diese Luft weniger, sodass der Temperaturunterschied zwischen Kühlrippen und Luft grösser bleibt, und die schnellere Strömung verbessert den Wärmeübergang. Allerdings erzeugt der Motor bei höherer Geschwindigkeit auch mehr Abwärme.
\end{abcliste}
